% Curs 10 - Pagina 10 (6-XII-1994).
% Fragment pentru % Curs 10 - Pagina 10 (6-XII-1994).
% Fragment pentru % Curs 10 - Pagina 10 (6-XII-1994).
% Fragment pentru % Curs 10 - Pagina 10 (6-XII-1994).
% Fragment pentru \input{courses/pages/curs10-pagina10-ro}.
% Fara preambul si fara intrerupere fortata de pagina.
% Continuarea aplicatiei despre grupuri de ordin pq de pe pagina 9.

\subsection*{Curs 10 - Pagina 10}

\noindent La fel,
\[
|\operatorname{Syl}_q(G)|=1.
\]

\noindent $\Longrightarrow P\trianglelefteq G$, $Q\trianglelefteq G$.

\medskip
\noindent $aba^{-1}b^{-1}\in P\cap Q$.

\medskip
\noindent $P=\langle a\rangle\trianglelefteq G$
\[
\Longrightarrow ba^{-1}b^{-1}\in P
\Longrightarrow aba^{-1}b^{-1}\in P.
\]

\noindent $ab=c\quad\Longrightarrow\quad\operatorname{ord}(c)=pq$.

% Nota editoriala: manuscrisul nu scrie explicit P\cap Q={e},
% nici consecinta ab=ba. Ele rezulta din normalitatea lui P si Q,
% din |P|=p, |Q|=q si p\ne q, justificand ord(ab)=pq.
% Nu au fost introduse completari in textul transcris.
.
% Fara preambul si fara intrerupere fortata de pagina.
% Continuarea aplicatiei despre grupuri de ordin pq de pe pagina 9.

\subsection*{Curs 10 - Pagina 10}

\noindent La fel,
\[
|\operatorname{Syl}_q(G)|=1.
\]

\noindent $\Longrightarrow P\trianglelefteq G$, $Q\trianglelefteq G$.

\medskip
\noindent $aba^{-1}b^{-1}\in P\cap Q$.

\medskip
\noindent $P=\langle a\rangle\trianglelefteq G$
\[
\Longrightarrow ba^{-1}b^{-1}\in P
\Longrightarrow aba^{-1}b^{-1}\in P.
\]

\noindent $ab=c\quad\Longrightarrow\quad\operatorname{ord}(c)=pq$.

% Nota editoriala: manuscrisul nu scrie explicit P\cap Q={e},
% nici consecinta ab=ba. Ele rezulta din normalitatea lui P si Q,
% din |P|=p, |Q|=q si p\ne q, justificand ord(ab)=pq.
% Nu au fost introduse completari in textul transcris.
.
% Fara preambul si fara intrerupere fortata de pagina.
% Continuarea aplicatiei despre grupuri de ordin pq de pe pagina 9.

\subsection*{Curs 10 - Pagina 10}

\noindent La fel,
\[
|\operatorname{Syl}_q(G)|=1.
\]

\noindent $\Longrightarrow P\trianglelefteq G$, $Q\trianglelefteq G$.

\medskip
\noindent $aba^{-1}b^{-1}\in P\cap Q$.

\medskip
\noindent $P=\langle a\rangle\trianglelefteq G$
\[
\Longrightarrow ba^{-1}b^{-1}\in P
\Longrightarrow aba^{-1}b^{-1}\in P.
\]

\noindent $ab=c\quad\Longrightarrow\quad\operatorname{ord}(c)=pq$.

% Nota editoriala: manuscrisul nu scrie explicit P\cap Q={e},
% nici consecinta ab=ba. Ele rezulta din normalitatea lui P si Q,
% din |P|=p, |Q|=q si p\ne q, justificand ord(ab)=pq.
% Nu au fost introduse completari in textul transcris.
.
% Fara preambul si fara intrerupere fortata de pagina.
% Continuarea aplicatiei despre grupuri de ordin pq de pe pagina 9.

\subsection*{Curs 10 - Pagina 10}

\noindent La fel,
\[
|\operatorname{Syl}_q(G)|=1.
\]

\noindent $\Longrightarrow P\trianglelefteq G$, $Q\trianglelefteq G$.

\medskip
\noindent $aba^{-1}b^{-1}\in P\cap Q$.

\medskip
\noindent $P=\langle a\rangle\trianglelefteq G$
\[
\Longrightarrow ba^{-1}b^{-1}\in P
\Longrightarrow aba^{-1}b^{-1}\in P.
\]

\noindent $ab=c\quad\Longrightarrow\quad\operatorname{ord}(c)=pq$.

% Nota editoriala: manuscrisul nu scrie explicit P\cap Q={e},
% nici consecinta ab=ba. Ele rezulta din normalitatea lui P si Q,
% din |P|=p, |Q|=q si p\ne q, justificand ord(ab)=pq.
% Nu au fost introduse completari in textul transcris.
